% GENERATED FILE. Edit canonical lessons, then run npm run generate:books.
% canonical-source-hash: fnv1a64:a03f463835757e52
% canonical-lessons: GE-C282-ueberfallen, GE-C282-ueberfall, GE-C282-rauben, GE-C282-bedrohen, GE-C282-waffe

\chapter{Der Überfall}
\label{ch:ge-der-uberfall}
\hlchaptermodality{282}

\begin{chapteropening}
\textbf{By the end of this chapter:} \emph{I can report a robbery: to hold up, a hold-up, to rob, to threaten and a weapon.}

Complete the last lesson of chapter 282: I can report a robbery: to hold up, a hold-up, to rob, to threaten and a weapon.
\end{chapteropening}

\section[überfallen]{überfallen — to hold up, to attack}

\label{lesson:GE-C282-ueberfallen}



\begin{quote}
Before the new one: say the German for to pinch, then the German for the loot.
\end{quote}

\subsection*{You'll want to know: überfallen}
\textbf{überfallen} — \textquotedblleft{}to hold up, to attack\textquotedblright{}.

\begin{quote}
Zwei Männer überfielen eine Tankstelle.
\end{quote}

\begin{quote}
\emph{Two men held up a petrol station.}
\end{quote}

Not separable: \emph{sie überfielen}, \emph{sie haben überfallen}.

\begin{cousinweb}
\emph{überfallen} is inherited from Middle High German \emph{übervallen} (EtymDB), built on \emph{über} \textquotedblleft{}over\textquotedblright{} and \emph{fallen} \textquotedblleft{}to fall\textquotedblright{}: to fall upon someone.
\end{cousinweb}

\begin{grammarlens}[title={What you've built}]
One more word for telling, explaining and describing at B1.
\end{grammarlens}

\subsection*{Your turn}
\begin{itemize}
\raggedright
  \item \emph{Say it:} \emph{überfallen}
  \item \emph{Say it:} \emph{überfallen}, once more
  \item \emph{Say it:} a sentence of your own with \emph{überfallen}
  \item \emph{Recall:} say the German for to pinch, then the German for the loot, then say \emph{überfallen} again
\end{itemize}

\begin{tcolorbox}[breakable,colback=teal!4,colframe=teal!35!black,title={Before you move on}]
What does überfallen mean? (\textquotedblleft{}To hold up, to attack\textquotedblright{}.) Say it once more.
\end{tcolorbox}
\section[der Überfall]{der Überfall — the hold-up, the raid}

\label{lesson:GE-C282-ueberfall}



\begin{quote}
Before the new one: say the German for the loot, then the German for to hold up.
\end{quote}

\subsection*{You'll want to know: der Überfall}
\textbf{der Überfall} — \textquotedblleft{}the hold-up, the raid\textquotedblright{}.

\begin{quote}
Bei dem Überfall wurde niemand verletzt.
\end{quote}

\begin{quote}
\emph{Nobody was hurt in the hold-up.}
\end{quote}

Plural: \emph{die Überfälle}.

\begin{cousinweb}
\emph{Überfall} is built on \emph{überfallen}, Middle High German \emph{übervallen} \textquotedblleft{}to fall upon\textquotedblright{} (EtymDB), from \emph{fallen} \textquotedblleft{}to fall\textquotedblright{}.
\end{cousinweb}

\begin{grammarlens}[title={What you've built}]
One more word for telling, explaining and describing at B1.
\end{grammarlens}

\subsection*{Your turn}
\begin{itemize}
\raggedright
  \item \emph{Say it:} \emph{der Überfall}
  \item \emph{Say it:} \emph{der Überfall}, once more
  \item \emph{Say it:} a sentence of your own with \emph{der Überfall}
  \item \emph{Recall:} say the German for the loot, then the German for to hold up, then say \emph{der Überfall} again
\end{itemize}

\begin{tcolorbox}[breakable,colback=teal!4,colframe=teal!35!black,title={Before you move on}]
What does der Überfall mean? (\textquotedblleft{}The hold-up, the raid\textquotedblright{}.) Say it once more.
\end{tcolorbox}
\section[rauben]{rauben — to rob, to steal by force}

\label{lesson:GE-C282-rauben}



\begin{quote}
Before the new one: say the German for to hold up, then the German for the hold-up.
\end{quote}

\subsection*{You'll want to know: rauben}
\textbf{rauben} — \textquotedblleft{}to rob, to steal by force\textquotedblright{}.

\begin{quote}
Der Mann raubte einer Frau die Tasche.
\end{quote}

\begin{quote}
\emph{The man robbed a woman of her bag.}
\end{quote}

\emph{Jemandem etwas rauben}: the person goes in the dative. \emph{Der Raub} — the robbery.

\begin{cousinweb}
\emph{rauben} is inherited from Old High German \emph{roubōn}, which EtymDB traces to Proto-Germanic \emph{raubōną} and further to a root meaning \textquotedblleft{}to snatch\textquotedblright{}.
\end{cousinweb}

\begin{grammarlens}[title={What you've built}]
One more word for telling, explaining and describing at B1.
\end{grammarlens}

\subsection*{Your turn}
\begin{itemize}
\raggedright
  \item \emph{Say it:} \emph{rauben}
  \item \emph{Say it:} \emph{rauben}, once more
  \item \emph{Say it:} a sentence of your own with \emph{rauben}
  \item \emph{Recall:} say the German for to hold up, then the German for the hold-up, then say \emph{rauben} again
\end{itemize}

\begin{tcolorbox}[breakable,colback=teal!4,colframe=teal!35!black,title={Before you move on}]
What does rauben mean? (\textquotedblleft{}To rob, to steal by force\textquotedblright{}.) Say it once more.
\end{tcolorbox}
\section[bedrohen]{bedrohen — to threaten}

\label{lesson:GE-C282-bedrohen}



\begin{quote}
Before the new one: say the German for the hold-up, then the German for to rob.
\end{quote}

\subsection*{You'll want to know: bedrohen}
\textbf{bedrohen} — \textquotedblleft{}to threaten\textquotedblright{}.

\begin{quote}
Er bedrohte den Verkäufer mit einem Messer.
\end{quote}

\begin{quote}
\emph{He threatened the shop assistant with a knife.}
\end{quote}

\emph{Jemanden mit etwas bedrohen}. The plain verb \emph{drohen} takes the dative: \emph{er drohte ihm}.

\begin{cousinweb}
\emph{bedrohen} is built on \emph{drohen} \textquotedblleft{}to threaten\textquotedblright{}, from Old High German \emph{drouwen} (EtymDB), with the prefix \emph{be-}.
\end{cousinweb}

\begin{grammarlens}[title={What you've built}]
One more word for telling, explaining and describing at B1.
\end{grammarlens}

\subsection*{Your turn}
\begin{itemize}
\raggedright
  \item \emph{Say it:} \emph{bedrohen}
  \item \emph{Say it:} \emph{bedrohen}, once more
  \item \emph{Say it:} a sentence of your own with \emph{bedrohen}
  \item \emph{Recall:} say the German for the hold-up, then the German for to rob, then say \emph{bedrohen} again
\end{itemize}

\begin{tcolorbox}[breakable,colback=teal!4,colframe=teal!35!black,title={Before you move on}]
What does bedrohen mean? (\textquotedblleft{}To threaten\textquotedblright{}.) Say it once more.
\end{tcolorbox}
\section[die Waffe]{die Waffe — the weapon}

\label{lesson:GE-C282-waffe}



\begin{quote}
Before the new one: say the German for to rob, then the German for to threaten.
\end{quote}

\subsection*{You'll want to know: die Waffe}
\textbf{die Waffe} — \textquotedblleft{}the weapon\textquotedblright{}.

\begin{quote}
Die Polizei fand keine Waffe.
\end{quote}

\begin{quote}
\emph{The police found no weapon.}
\end{quote}

Plural: \emph{die Waffen}.

\begin{cousinweb}
\emph{Waffe} is inherited from Old High German \emph{wāfan} \textquotedblleft{}weapon\textquotedblright{}, from Proto-Germanic \emph{wēpną} (EtymDB), the same source as English \emph{weapon}.
\end{cousinweb}

\begin{grammarlens}[title={What you've built}]
One more word for telling, explaining and describing at B1.
\end{grammarlens}

\subsection*{Your turn}
\begin{itemize}
\raggedright
  \item \emph{Say it:} \emph{die Waffe}
  \item \emph{Say it:} \emph{die Waffe}, once more
  \item \emph{Say it:} a sentence of your own with \emph{die Waffe}
  \item \emph{Recall:} say the German for to rob, then the German for to threaten, then say \emph{die Waffe} again
\end{itemize}

\begin{tcolorbox}[breakable,colback=teal!4,colframe=teal!35!black,title={Before you move on}]
What does die Waffe mean? (\textquotedblleft{}The weapon\textquotedblright{}.) Say it once more.
\end{tcolorbox}
